% XeLaTeX + ctex；在 docs/latent-model/latex/ 下编译：
%   latexmk -xelatex -interaction=nonstopmode -halt-on-error latent-model.tex

\documentclass[UTF8,a4paper,11pt]{ctexart}

\usepackage{amsmath,amssymb}
\usepackage{graphicx}
\usepackage{booktabs}
\usepackage{hyperref}
\usepackage{geometry}
\usepackage[backend=biber,style=numeric,sorting=none]{biblatex}
\addbibresource{refs.bib}

\geometry{margin=2.5cm}
\hypersetup{colorlinks=true, linkcolor=blue, citecolor=blue, urlcolor=blue}

\title{bdelf latent\_t5 / latent\_vae：设计规格与实现说明}
\author{bdelf 模型文档}
\date{\today}

\begin{document}
\maketitle

\begin{abstract}
本文档完整说明 bdelf 仓库中两个 Stage-1 连续 latent 模型族 \texttt{latent\_t5} 与 \texttt{latent\_vae} 的动机、结构、读出变体、损失函数与训练接口。
二者共享 T5-small 级张量形状（$E{=}512$，$d_{\mathrm{kv}}{=}64$，$d_{\mathrm{ff}}{=}2048$，encoder/decoder 各 6 层），但 decoder 拓扑与辅助损失不同。
本文档亦说明\textbf{为何}在 Cola 路线之外做这组对照，以及二者作为扩散语言模型（DLM）Stage-2 \emph{语义空间}的初步优劣分析。
工程路径、CLI 与配置表见同目录上级 \texttt{README.md}。
\end{abstract}

\section{动机与定位}

\subsection{连续 latent 语言模型}

变分自编码器（VAE）将观测 $x$ 与连续潜变量 $z$ 通过近似后验 $q_\phi(z|x)$ 与生成模型 $p_\theta(x|z)$ 联结~\cite{kingma2014auto}。
Cola DLM~\cite{cola2026dlm} 在 Stage-1 训练 Text VAE，损失为重建、KL 正则与 BERT 式 mask 辅助项，Stage-2 再在 latent 序列上训练扩散先验。

本仓库已有官方风格的 \texttt{cola\_vae}（384 维 SwiGLU 块因果 VAE）。
\texttt{latent\_t5} 与 \texttt{latent\_vae} 是\textbf{并列实验族}：在\textbf{相同 T5-small 维数}下对比两种 decoder 归纳偏置（自回归 + cross-attn \emph{vs.} 并行统一生成），并共享 encoder 实现以便控制变量。

\subsection{与仓库内其它模型的关系}

\begin{itemize}
  \item \textbf{非} HuggingFace 冻结 \texttt{t5-small}（ELF 家族用法）；内部为手写 Transformer。
  \item \textbf{非} \texttt{cola\_vae} 替换品；Cola Stage-2 仍只加载 \texttt{cola\_vae}。
  \item 二者注册为 \texttt，训练日志字段与 \texttt{cola\_vae} 同构（\texttt{recon\_ce}、\texttt{kl}、\texttt{mask} 等）。
\end{itemize}

\section{为何做这些对照实验}
\label{sec:why}

\subsection{DLM 两阶段视角}

连续 latent DLM（以 Cola 为代表~\cite{cola2026dlm}）将文本生成拆为两阶段：
\begin{enumerate}
  \item \textbf{Stage-1（本文档模型）}：学习 $q_\phi(z_{1:L}\mid x_{1:L})$ 与 $p_\theta(x_{1:L}\mid z_{1:L})$，得到逐 token 的连续表示 $z_t\in\mathbb{R}^{d_z}$；
  \item \textbf{Stage-2（DiT 先验）}：在 $z_{1:L}$ 上训练扩散/流模型 $p_\psi(z)$，推理时先采样 latent 再 decode 为文本。
\end{enumerate}
Stage-2 默认假设：$z$ 构成一个\textbf{平滑、可扩散、语义充分}的序列空间。
Stage-1 的 encoder 视野、瓶颈宽度、decoder 因子分解与辅助损失，都会改变 $z$ 的几何与可重建性，因而影响 DLM 上限。

\subsection{相对 \texttt{cola\_vae} 要回答的问题}

\texttt{cola\_vae} 已验证「块因果 encoder + 并行 decoder + BERT-mask」在 Cola 官方设定下可行，但其规模（384 维、4+4 层、SwiGLU 块）与 DiT 默认接口已固定。
本对照\textbf{不替换} Stage-2 现网路径，而是在\textbf{可控变量}下追问：

\begin{itemize}
  \item \textbf{Encoder 视野}：\texttt{latent\_t5} 可选双向或单向；\texttt{latent\_vae} 仅因果/块因果（禁止双向）。$z_t$ 是否更应携带全句上下文？
  \item \textbf{Decoder 归纳偏置}：自回归 + cross-attn \emph{vs.} 并行统一生成，何者更利于学习「可被扩散补全」的 $z$？
  \item \textbf{Latent 维度与读出}：$z\in\mathbb{R}^E$ \emph{vs.} $z\in\mathbb{R}^B$（$B{\ll}E$），扩散先验在低维是否更省算力、是否损失语义？
  \item \textbf{辅助损失形态}：span corruption（T5）\emph{vs.} BERT-mask（VAE），何者更能抑制 encoder 塌缩、同时保持 $z$ 对未见 token 的可预测性？
  \item \textbf{骨干对齐}：在 T5-small 张量形状下复用同一 \texttt{encdec}，排除「层数/宽度不同」带来的混杂，使对比聚焦在\textbf{拓扑与损失}。
\end{itemize}

\subsection{实验产出与判据（Stage-1）}

本仓库 Stage-1 训练可直接观测：
重建 CE（\texttt{recon\_ce}）、KL、辅助项（\texttt{mask}）、以及 held-out 重建。
后续若接入 Stage-2，建议以 \textbf{latent 扩散 NLL / 生成 PPL / 前缀条件一致性} 为语义空间的外部判据；
当前文档仅给出\textbf{结构层面的先验分析}（下一节），不预设哪条路线必然更优。

\section{作为 DLM 语义空间的优劣分析}
\label{sec:dlm-space}

以下从「Stage-2 扩散先验 + decode」需求出发，对两条路线作\textbf{定性}对比；$z$ 的维数记为 $d_z\in\{B,E\}$。

\subsection{共同前提}

二者均用 Cola 式 ELBO 训练，$z_{1:L}$ 为\textbf{逐位置}连续向量序列（非整句单向量）。
扩散模型通常在 $L\times d_z$ 张量上操作；故 $d_z$ 与 $L$ 共同决定先验难度。
默认 $B{=}64$ 时，\texttt{latent\_vae} 与 \texttt{latent\_t5}（\texttt{readout=b}）的 KL 均在 $B$ 维；
\texttt{readout=e} 则在 $E{=}512$ 维施加 KL，先验负担显著更重。

\subsection{\texttt{latent\_vae}：并行解码、低维 $z$}

\paragraph{潜在优势}
\begin{itemize}
  \item \textbf{与 Cola 范式一致}：并行 $p(x|z)$ 与「一次生成整句 latent 再 decode」的推理图一致；Stage-2 采样 $z_{1:L}$ 后单次前向即可得全文，延迟低。
  \item \textbf{低维 $z\in\mathbb{R}^B$}：扩散状态空间为 $L{\times}B$；相同网络宽度下，先验参数量与去噪步计算通常低于 $L{\times}E$。
  \item \textbf{BERT-mask 辅助}：迫使 encoder 在输入被遮挡时仍预测语义，减轻 $z$ 仅编码「表面共现」、decoder 死记 $x$ 的塌缩~\cite{devlin2019bert}。
  \item \textbf{与块扩散/块内双向先验兼容}：若 Stage-2 采用块级噪声（Cola 中 block 与 DiT 对齐），可选 \texttt{block\_size>1} 的 encoder 与块内双向 $z$ 更同构（默认 \texttt{block\_size=1} 为因果基线）。
\end{itemize}

\paragraph{潜在劣势}
\begin{itemize}
  \item \textbf{因果 encoder（默认）}：$z_t$ 仅依赖 $x_{\le t}$，不含后文语境；对需要全局一致性的主题/指代，$z$ 可能欠充分，需 decoder 从 $z_{1:t}$ 外推补全。
  \item \textbf{并行 decoder 与 AR 评测错位}：训练目标为「给定全长 $z$ 同时预测各 token」，与自回归 PPL 度量的因子分解不一致；生成质量需用重建或专用指标。
  \item \textbf{信息瓶颈}：$B{=}64$ 仍可能限制极细粒度语义；若 Stage-2 需更细风格/事实，可能表现为重建尚可但扩散样本偏模糊。
\end{itemize}

\subsection{\texttt{latent\_t5}：encoder 与 decoder self-attn 同模式}

\paragraph{潜在优势}
\begin{itemize}
  \item \textbf{双向（默认 \texttt{bidirectional=true}）}：encoder 与 decoder self-attn 均见全句；$z_t$ 可编码左右语境；decoder 从 $z$ 并行重建。作为 DLM「语义坐标」更贴近 BERT/T5 encoder 表征~\cite{devlin2019bert,raffel2020t5}。
  \item \textbf{单向消融}：\texttt{bidirectional=false} 时 encoder \textbf{与} decoder self-attn 均为逐 token 因果，AR + cross-attn，可与 \texttt{latent\_vae} 的因果视野对照。
  \item \textbf{单向 AR 重建}：$p(x|z)$ 为自回归，与主流 LM 评测兼容；前缀条件时 decoder 自然续写。
  \item \textbf{span 辅助}：腐蚀 encoder 输入、decoder 仍重建原句，类似去噪自编码，可能使 $z$ 对局部缺失更鲁棒，利于 Stage-2 对 $z$ 的随机扰动。
  \item \textbf{\texttt{readout=b}}：与 VAE 共享 $z\in\mathbb{R}^B$，便于在\textbf{同一扩散维度}下比较「双向并行」与「因果+并行 / 因果 AR」谁学得更好。
\end{itemize}

\paragraph{潜在劣势}
\begin{itemize}
  \item \textbf{\texttt{readout=e} 的高维先验}：$z\in\mathbb{R}^E$ 使 Stage-2 扩散在 $L{\times}512$ 空间进行，训练与采样成本高，且高维高斯先验与 $\mathcal{N}(0,I)$ 的 KL 更易与重建拉扯。
  \item \textbf{AR decoder 与并行扩散的因子分解不一致}：Stage-2 通常并行去噪各 $z_t$，而 Stage-1 用 AR 学 $p(x|z)$；若 $z$ 各维强依赖 AR 解码顺序，扩散「同步更新 $z$」可能引入 train--test 间隙。
  \item \textbf{无条件生成}：宽先验下从 $\mathcal{N}(0,I)$ 采 $z$ 再 AR，质量往往弱于「encode 前缀」；DLM 若强调无条件文本生成，需依赖 Stage-2 学好 $p(z)$。
  \item \textbf{算力}：cross-attn decoder 每层较 VAE 并行 decoder 多一组 cross-attn，Stage-1 训练更贵。
\end{itemize}

\subsection{对照小结}

\begin{center}
\small
\begin{tabular}{p{2.2cm}p{4.8cm}p{4.8cm}}
\toprule
维度 & \texttt{latent\_vae} & \texttt{latent\_t5} \\
\midrule
$z$ 默认维数 & $B{=}64$ & $E{=}512$ 或 $B{=}64$（\texttt{readout}） \\
Encoder 语境 & 因果/块因果（禁止双向） & 双向（默认）或单向 \\
Decoder & 并行，一次出全长 & AR + cross-attn \\
辅助损失 & BERT-mask & span corruption \\
Stage-2 先验难度 & 低（$L{\times}B$） & \texttt{e} 高；\texttt{b} 与 VAE 同级 \\
推理图 & 采 $z$ $\to$ 并行 decode & 采/encode $z$ $\to$ AR \\
与 Cola 现网 & 拓扑不同；非即插 & 同左 \\
\bottomrule
\end{tabular}
\end{center}

\paragraph{实践建议（假设性）}
若 Stage-2 预算有限、优先低维扩散与单次 decode，可优先关注 \texttt{latent\_vae} 及 \texttt{latent\_t5}（\texttt{readout=b}）。
若更看重 $z_t$ 的全句语义与 AR 下游兼容，可优先 \texttt{latent\_t5}（\texttt{readout=e} 或 \texttt{b}）。
最终应以 Stage-1 重建、KL--重建折中、以及接入 DiT 后的生成指标为准；本文档不提供实验结论。

\section{共享骨干：形状与算子}

\subsection{张量形状与默认超参（\texttt{100m.yaml}）}
\label{sec:shape}

\begin{center}
\begin{tabular}{lll}
\toprule
键 & \texttt{latent\_t5} & \texttt{latent\_vae} \\
\midrule
\texttt{n\_embd} / \texttt{d\_kv} / \texttt{d\_ff} / 层数 & \multicolumn{2}{c}{512 / 64 / 2048 / enc6+dec6} \\
\texttt{latent\_dim}（$B$） & \multicolumn{2}{c}{16} \\
\texttt{max\_seq\_len} & \multicolumn{2}{c}{4096} \\
\texttt{tokenizer} & \multicolumn{2}{c}{GPT-2（OWT 管线）} \\
encoder / decoder self-attn & \texttt{bidirectional}（默认 true：二者双向并行重建；false：二者因果 AR） & 禁止双向；\texttt{block\_size}（默认 1） \\
读出 & \texttt{readout=e|b} & 固定 $B$ 读出 \\
辅助损失 & \texttt{lambda\_span} + span 腐蚀 & \texttt{lambda\_mask} + BERT-mask \\
\texttt{beta\_kl} & \multicolumn{2}{c}{0.1} \\
\bottomrule
\end{tabular}
\end{center}

与 T5-small 的 $d_{\mathrm{model}}, d_{\mathrm{kv}}, d_{\mathrm{ff}}, 6{+}6$ 层对齐；\textbf{未}复现 T5 相对位置偏置、RMSNorm 与 ReLU FFN，而采用本仓库 AR 栈惯例：pre-norm LayerNorm、RoPE~\cite{su2024rope}、GELU FFN。

\subsection{共享模块 \texttt{models/latent/encdec/}}

\begin{description}
  \item[\texttt{SelfAttention}] $Q,K,V$ 投影至 $H \cdot d_{\mathrm{kv}}$；RoPE 施加于 $Q,K$。
  \item[\texttt{TransformerBlock}] pre-norm self-attn + FFN。
  \item[\texttt{LatentEncoder}] token embedding + $N_{\mathrm{enc}}$ 层 block；支持三种注意力模式（第~\ref{sec:attn} 节）。
  \item[\texttt{PosteriorBReadout}] $\mu,\log\sigma^2$ 均为 $\mathbb{R}^{E}\!\to\!\mathbb{R}^{B}$；$\mu$ 经无仿射 LayerNorm。
  \item[\texttt{PosteriorEReadout}] 瓶颈 $\mathbb{R}^{E}\!\to\!\mathbb{R}^{B}$，再 $\mathbb{R}^{B}\!\to\!\mathbb{R}^{E}$ 得 $\mu,\log\sigma^2$。
\end{description}

\section{Encoder 注意力三路}
\label{sec:attn}

设 encoder 隐状态为 $h \in \mathbb{R}^{L \times E}$。

\begin{enumerate}
  \item \textbf{双向}（\texttt{bidirectional=true}，\texttt{latent\_t5} 默认）：encoder \textbf{与} decoder self-attn 均无因果 mask。Decoder 从 $z$ 并行重建（不 teacher-force $x$），以免双向 self-attn 看见未来 token。\textbf{仅} \texttt{latent\_t5} 可配置。
  \item \textbf{逐 token 因果}（\texttt{bidirectional=false}, \texttt{block\_size=1}）：encoder 与 decoder self-attn 均为 \texttt{is\_causal=True}；decoder 仍 AR + cross-attn。\texttt{latent\_t5} 可通过 \texttt{--set model.bidirectional=false} 选用；\texttt{latent\_vae} 固定此模式（默认 \texttt{block\_size=1}）。
  \item \textbf{块因果}（\texttt{block\_size>1}，可选）：块内双向、块间因果；mask 与 \texttt{cola\_vae} 共用 \texttt{block\_causal\_mask\_mod}。要求 $L \bmod \texttt{block\_size} = 0$。\texttt{latent\_vae} 默认 \texttt{block\_size=1}，即不启用块注意力。
\end{enumerate}

\texttt{latent\_vae} \textbf{禁止} \texttt{bidirectional=true}（配置层与实现层均拒绝；encoder 写死 \texttt{False}）。

\section{瓶颈与读出}

\subsection{记号}

后验参数化为对角高斯 $q(z|x)=\mathcal{N}(\mu,\mathrm{diag}(\sigma^2))$，网络输出 $\log\sigma^2$ 并 clamp 至 $[-30,20]$。
重参数化 $z=\mu+\sigma\odot\varepsilon$，$\varepsilon\sim\mathcal{N}(0,I)$。
KL 为
\begin{equation}
  \mathrm{KL}(q\|p_0) = -\tfrac{1}{2}\,\mathbb{E}\bigl[1 + \log\sigma^2 - \mu^2 - e^{\log\sigma^2}\bigr],
\end{equation}
对 $z$ 所在空间各维取平均（\texttt{kl\_gaussian}）。

\subsection{\texttt{latent\_vae}：$z \in \mathbb{R}^{L \times B}$}

对 encoder 输出 $h_{\mathrm{enc}}$：
\begin{align}
  \mu &= \mathrm{LN}\bigl(W_\mu h_{\mathrm{enc}}\bigr), \quad W_\mu: \mathbb{R}^{E}\!\to\!\mathbb{R}^{B}, \\
  \log\sigma^2 &= W_{\sigma} h_{\mathrm{enc}}, \quad W_{\sigma}: \mathbb{R}^{E}\!\to\!\mathbb{R}^{B}.
\end{align}
Decoder 输入为 $\mathrm{from\_latent}(z)$，$W_z: \mathbb{R}^{B}\!\to\!\mathbb{R}^{E}$，再经与 encoder 同型的 self-attn 栈（默认同为逐 token 因果），一次前向得到全长 logits（\emph{非}自回归）。

\subsection{\texttt{latent\_t5}：两种读出}

\subsubsection{\texttt{readout=e}（默认）：$z \in \mathbb{R}^{L \times E}$}

\begin{align}
  b &= W_{E\to B}\, h_{\mathrm{enc}}, \\
  \mu &= \mathrm{LN}\bigl(W_{B\to E}\, b\bigr), \quad
  \log\sigma^2 = W'_{B\to E}\, b.
\end{align}
瓶颈 $b\in\mathbb{R}^B$ 为中间压缩；随机变量 $z$ 仍在 $E$ 维，作为 decoder \textbf{cross-attention} 的 memory。
Cross-attn 的 $K,V$ 由 $z$ 经 $\mathbb{R}^{E}\!\to\!H\cdot d_{\mathrm{kv}}$ 线性投影（标准 T5 decoder 在 $E$ 维 memory 上的情形）。

\subsubsection{\texttt{readout=b}：$z \in \mathbb{R}^{L \times B}$}

$\mu,\log\sigma^2$ 与 \texttt{latent\_vae} 相同（均在 $B$ 维）。
\textbf{不}先将 $z$ 扩维到 $E$ 再作 memory；decoder 每层 cross-attn 的 $K,V$ 为
\begin{equation}
  K = W_K z,\quad V = W_V z,\quad W_K,W_V: \mathbb{R}^{B}\!\to\!H\cdot d_{\mathrm{kv}},
\end{equation}
而 $Q$ 仍来自 decoder 隐状态 $h_{\mathrm{dec}}\in\mathbb{R}^E$。
Self-attention 的 $Q,K,V$ 始终在 $E$ 维。

此变体使 T5 路径上的 KL 与 VAE 一样在 $B{=}16$ 维平均，便于公平对比瓶颈宽度；代价是 cross-attn memory 信息带宽更低。

\section{Decoder 与生成}

\subsection{\texttt{latent\_t5}：cross-attn decoder（self-attn 与 encoder 同模式）}

每层：self-attn（RoPE；\texttt{bidirectional} 与 encoder 一致）$\to$ cross-attn($h_{\mathrm{dec}}, z$) $\to$ FFN。

\paragraph{训练}
\begin{itemize}
  \item \texttt{bidirectional=true}：decoder 从 $z$ 起步（\texttt{readout=e} 时 $z$ 已是 $E$ 维；\texttt{readout=b} 时经 $B\to E$），一次预测全长 $x$。数据仍因 \texttt{full\_sequence\_training=True} 传入全长。
  \item \texttt{bidirectional=false}：构造 $\tilde{x} = [\mathrm{BOS}, x_1,\ldots,x_{L-1}]$ 做教师强制 AR。\textbf{不}使用 GPT 式外层 \texttt{batch[:,:-1]} 移位（否则 encoder 会丢失末 token）。
\end{itemize}

\paragraph{生成}
\begin{itemize}
  \item 双向：一次 decode 全长；有前缀则 encode（前缀+$L$ pad）后并行重建并覆盖前缀。
  \item 单向：有前缀则 encode 得 $z$ 再 AR 续写；无条件从 $\mathcal{N}(0,I)$ 采 $z$ 再 AR。
\end{itemize}

\subsection{\texttt{latent\_vae}：统一生成}

Decoder 仅接收 $\mathrm{from\_latent}(z)$，与 Cola \texttt{decode\_logits} 同型。
无条件从先验采 $z\in\mathbb{R}^{L\times B}$ 再一次 decode；有前缀则 encode（前缀+$L$ pad）后并行 decode 并保留前缀。

\section{损失函数}

\subsection{总目标}

训练态统一为
\begin{equation}
  \mathcal{L} = \mathcal{L}_{\mathrm{recon}} + \beta\,\mathrm{KL}(q\| \mathcal{N}(0,I)) + \lambda\, \mathcal{L}_{\mathrm{aux}},
\end{equation}
默认 $\beta{=}0.1$，$\lambda{=}1$（与 \texttt{cola\_vae} 默认一致）。

\subsection{重建项 $\mathcal{L}_{\mathrm{recon}}$}

\begin{description}
  \item[\texttt{latent\_t5}] 双向（默认）：并行 CE $-\sum_t \log p_\theta(x_t \mid z)$。单向：教师强制 AR CE $-\sum_t \log p_\theta(x_t \mid x_{<t}, z)$。
  \item[\texttt{latent\_vae}] 并行 CE：$-\sum_t \log p_\theta(x_t \mid z)$，所有位置同时预测。
\end{description}

Held-out 评测走模型 \texttt{forward} 的重建 CE（T5 默认并行；单向为 teacher-forced AR；VAE 并行；\texttt{config/train/eval/latent/default.yaml}）。

\subsection{辅助项 $\mathcal{L}_{\mathrm{aux}}$}

\paragraph{\texttt{latent\_vae}：BERT-mask}
以概率 \texttt{mask\_ratio}（默认 0.15）将 token 换为专用 \texttt{mask\_token\_id}$=V$（$V$ 为词表大小）。
第二次前向：encode 被 mask 的序列，仅在 mask 位置对原始 $x$ 计 CE~\cite{devlin2019bert}。
固定张量形状（\texttt{ignore\_index} 掩掉非 mask 位），便于 \texttt{torch.compile}。

\paragraph{\texttt{latent\_t5}：原地 span CE}
受 T5 span corruption 启发~\cite{raffel2020t5}，但\textbf{不}采用变长 decoder 目标序列（与定长训练图不兼容）。
流程：
\begin{enumerate}
  \item 按 span 级 mask 覆盖约 15\% token（平均 span 长 3）；
  \item encoder 输入将 span 换为 sentinel id $V{+}k$，$k\in\{0,\ldots,99\}$；
  \item decoder 仍 AR 全长原始 $x$；
  \item CE 仅在腐蚀位置计算。
\end{enumerate}
100 个 sentinel 仅扩展 encoder embedding，\textbf{不}进入 \texttt{lm\_head}。

\section{训练系统接口}

\subsection{CLI 与配置哈希}

\begin{verbatim}
train.py --model {latent_t5|latent_vae} --config 100m-{fast|full} \
  --dataset owt --preprocess default --generate eval
\end{verbatim}

Checkpoint：\texttt{cache/checkpoints/\{fast|full\}/latent/\{model\}/\{hash\}/}。
\texttt{readout}、\texttt{latent\_dim}、层数等进入 \texttt{config-hash}；\texttt{world\_size} 不进哈希。

\subsection{日志字段}

Latent 主表：\texttt{recon\_ce}、\texttt{kl}、\texttt{mask}、\texttt{token\_acc}、\texttt{mask\_acc}、\texttt{beta\_kl}、\texttt{lambda\_mask}（T5 的 $\lambda$ 写入 \texttt{lambda\_mask} 键，值为 \texttt{lambda\_span}）。

\subsection{优化器}

Muon 识别 \texttt{attn.qkv}、\texttt{attn.proj}、\texttt{mlp.c\_fc}、\texttt{mlp.c\_proj}；
读出层、\texttt{from\_latent}、\texttt{lm\_head}、embedding 走 AdamW。

\section{长度课程（Stage-1）}
\label{sec:curriculum}

本节为 \texttt{latent\_t5} / \texttt{latent\_vae} 的\textbf{长度课程}实验协议（非 Cola Stage-2 DiT）。
与 Cola 现网 Stage-2 无关；两条 latent 模型共用同一日程。
工程命令与 dataloader 接线见 \texttt{README.md}。

\subsection{预算与数据原则}

全程 \textbf{10B 有效 token}（pad 不计入科学预算）。
OWT 一整遍约 8--9B；10B 约 1.1 遍，多出的量用于长段外推与短桶回放。

\begin{itemize}
  \item 配比按\textbf{有效 token 抽样比例}，不是条数、不是 pad 后张量长度。
  \item \textbf{同一阶段不混} \texttt{owt-seg512} 与 \texttt{owt-bucket}：二者 512 切分规则不同，不是同一分布。
  \item 同一步只含同一桶。阶段内训练图长最多两档：短桶并入次长档（S1/S2 钉死 512；S3 为 $512$ 或 $1024$；S4 为 $1024$ 或 $2048$），微批按阶段最大 $L$ 固定。不按桶原生四档切换（\texttt{torch.compile} 多图与分配器碎片），也不把短序列垫到阶段最大 $L$（注意力 $L^2$）。Pad 不计损失。
\end{itemize}

\subsection{四阶段与配比}

每步图 token 在阶段最大 $L$ 时约 $262\text{K}$（$=$ 峰值 $L \times$ 全局条数）。短档步按该档 $L$ 计。

\begin{center}
\small
\begin{tabular}{lrrlrr}
\toprule
阶段 & 有效 token & 数据集 & 图 $L$ & 条数 & 每步 pad-token \\
\midrule
S1 & 3.0B & 仅 \texttt{owt-seg512} & 512 & 512 & $512^2$ \\
S2 & 1.5B & \texttt{owt-bucket} 256+512 桶 & 512 & 512 & $512^2$ \\
S3 & 2.5B & \texttt{owt-bucket} 256+512+1024 & 1024 & 256 & $256{\times}1024$ \\
S4 & 3.0B & \texttt{owt-bucket} 四桶 & 2048 & 128 & $128{\times}2048$ \\
\bottomrule
\end{tabular}
\end{center}

\begin{center}
\small
\begin{tabular}{lrrrrr}
\toprule
来源 & S1 & S2 & S3 & S4 & 合计 \\
\midrule
\texttt{owt-seg512} & 100\% & --- & --- & --- & 3.00B \\
bucket 256 & --- & 30\% & 10\% & 5\% & 0.85B \\
bucket 512 & --- & 70\% & 20\% & 10\% & 1.85B \\
bucket 1024 & --- & --- & 70\% & 20\% & 2.35B \\
bucket 2048 & --- & --- & --- & 65\% & 1.95B \\
\midrule
阶段合计 & 3.00B & 1.50B & 2.50B & 3.00B & \textbf{10.00B} \\
\bottomrule
\end{tabular}
\end{center}

S1 仅在定长 seg512 上训稳重建/KL/辅助三项；S2 换 bucket 短桶切分且\textbf{不改}优化器几何（仍 $512{\times}512$ 条）。
S3/S4 为长度外推：新长度桶占 70\%/65\%，短桶回放递减。

\subsection{优化器：一条 WSD，不按阶段重置}

\textbf{数据课程 $\neq$ 优化器课程。}
Stage-1 VAE 的 KL--重建折中由 Adam/Muon 动量维持；阶段边界重置优化器易在几步内打飞后验。
故全程一条 warmup--stable--decay（WSD）轨迹，EMA 不断档：

\begin{center}
\begin{tabular}{lrl}
\toprule
段 & 有效 token & 学习率 \\
\midrule
warmup（仅 S1 开头一次） & $\sim$150M & 线性 $0 \to 2{\times}10^{-3}$ \\
stable & $\sim$9.2B & 恒定 $2{\times}10^{-3}$ \\
decay（仅 S4 末） & 0.8B & cosine $\to 2{\times}10^{-4}$ \\
\bottomrule
\end{tabular}
\end{center}

Muon 与 AdamW 共用峰值 LR；$\beta_1{=}0.9$，$\beta_2{=}0.95$；无 weight decay；$\beta_{\mathrm{kl}}{=}0.1$，$\lambda{=}1$ 四阶段不变。
阶段边界\textbf{不}改 LR、\textbf{不}重置动量或 EMA。
若必须换 checkpoint 目录，须整包续训并按总 10B 的 token 进度续算 WSD，禁止本阶段 step$=0$ 再 warmup。

S3/S4 入口同时改 $L$ 与全局条数时，$|g|$ 尺度会短暂错配；靠 $\beta_2$ 的几十步自应消化，\textbf{不}靠重置或阶梯降 LR。

\subsection{阶段闸与观察窗}

\begin{itemize}
  \item \textbf{S1$\to$S2}：dev 重建 CE 走平；KL 不塌缩、不发散；辅助 \texttt{mask\_acc} 明显高于随机。
  \item \textbf{S2$\to$S3}：\texttt{owt-bucket} 的 512 桶重建不差于 S1 结束（勿用 seg512 dev 冒充）。
  \item \textbf{S3$\to$S4}：1024 桶重建已改善；256/512 桶无遗忘；S3 入口观察窗已过。
\end{itemize}

S3、S4 各在阶段开头留 \textbf{0.2B} 观察窗（计入该阶段预算）：允许 CE/KL 尖峰；窗内不降 LR、不重置优化器。
S4 的 cosine 退火仅在观察窗之后的最后 0.8B 进行。
评测须\textbf{按桶}报告 \texttt{recon\_ce}、KL、\texttt{mask\_acc}，并保留 bucket-512 遗忘对照；禁止单条 \texttt{eval\_loss} 决策。

\section{设计选择与局限}

\begin{enumerate}
  \item \textbf{RoPE 而非 T5 相对偏置}：与仓库 AR/块因果路径统一；双向 encoder 语义对齐 T5，算子不对齐 HF 权重。
  \item \textbf{GPT-2 词表}：与 OWT 预处理一致；未 tie 的 embedding 与 \texttt{lm\_head} 占模型参数较大头，高于 32k SPM 的 T5-small 计数。
  \item \textbf{KL 空间}：\texttt{readout=e} 时 KL 在 $E{=}512$；与 \texttt{latent\_vae} 的 $B{=}16$ 后验容量不对等；\texttt{readout=b} 可缓解。
  \item \textbf{span 辅助损失}：固定长度「原地 CE」是工程折中，非严格 T5 预训练目标。
  \item \textbf{无条件生成}：latent AR 依赖宽先验 $z$；实用场景建议前缀 encode。
\end{enumerate}

\section{小结}

\texttt{latent\_t5} 与 \texttt{latent\_vae} 在共享 T5-small 维数 encoder 与 Cola 式 ELBO 下，分别检验「双向并行重建（默认可切单向 AR）+ span 辅助」与「因果 + 并行解码 + BERT-mask」两条 Stage-1 路线（动机与 DLM 语义空间讨论见第~\ref{sec:why}--\ref{sec:dlm-space} 节）。
读出变体 \texttt{readout=e|b} 控制随机变量所在空间与 cross-attn memory 维数，便于消融瓶颈与 KL 尺度。
长度课程（10B、四阶段数据配比、WSD 优化器）见第~\ref{sec:curriculum} 节。
工程细节以代码与 \texttt{README.md} 为准；本文档描述设计意图与完整行为规格。

\printbibliography

\end{document}
